\documentclass[10pt,a4paper]{article}
\usepackage[margin=21mm]{geometry}
\usepackage{amsmath,amssymb,booktabs,array,longtable,xcolor,hyperref}
\hypersetup{colorlinks=true,linkcolor=blue,urlcolor=blue}
\setlength{\parskip}{5pt}
\setlength{\parindent}{0pt}
\newcommand{\code}[1]{\texttt{#1}}
\title{Anchor Selection in \texttt{mp\_isdf}\\\large Which Algorithms Change the Sketch?}
\author{Codebase guide \quad (\texttt{cpp/src/anchor\_select*.cpp}, \texttt{python/isdf.py})}
\date{Generated from the current source tree}
\begin{document}
\maketitle
\begin{abstract}
Anchor indices are grid sites. This guide groups every selector exposed by
\code{python/isdf.py} by its actual sketch matrix, rather than by whether it
uses randomness. In particular, randomized candidate sampling and randomized
initialization are \emph{not} Gaussian sketching. A new weighting or
normalization changes the matrix being sketched, but does not automatically
constitute a new random embedding distribution.
\end{abstract}

\section{Notation and decision map}
Let $g=1,\ldots,n$ index eligible grid sites, $w_g$ be the signed quadrature
weight, $a_g\in\mathbb R^d$ the AO-value vector, and $f_g\in\mathbb R^d$
the requested concatenated and scaled feature vector (AO, derivative, or
occupied-orbital channels). In the general-feature mode, each axis channel
is normalized by its Frobenius norm over eligible sites and may receive an
axis weight. Let $p_g=a_g\otimes a_g$ be the ordered AO-pair column and
$q_g=\|a_g\|_2^2$, so $\|p_g\|_2=q_g$. The sketch has $\ell$ rows and
$n$ columns; $u_t,v_t$ are independent standard Gaussian $d$-vectors, shared
across grid batches. All the Khatri--Rao modes use the \emph{same} factorized
Gaussian map $S_{t,(\mu,\nu)}=u_{t\mu}v_{t\nu}/\sqrt\ell$ and apply QR with
column pivoting (QRCP) to its sketched columns. They do not form the full
ordered pair matrix. Differences below concern $f_g$, column weights, and
normalization, not the random map itself.

\small
\begin{longtable}{@{}p{0.28\linewidth}p{0.29\linewidth}p{0.34\linewidth}@{}}
\toprule
CLI selector / technique & Sketch / random mechanism & What changes? \\
\midrule\endhead
\code{qrcp\_dense} & None: QRCP of full weighted feature pairs & Baseline, direct deterministic selection. \\
\code{qrcp\_online} & None: implicit full weighted feature pairs & Algorithmic, bounded-memory multipass direct QR selection; no new embedding. \\
\code{quadrature} & Gaussian Khatri--Rao pair sketch & Any selected axes; quadrature-weighted pairs. \\
\code{density} & Same Gaussian Khatri--Rao map & Replace quadrature column weight with importance; non-AO axes NYI. \\
\code{density\_col\_norm} & Same Gaussian Khatri--Rao map & Importance plus pair-column normalization; non-AO axes NYI. \\
\code{rbrp} & No Gaussian sketch; residual-weighted candidate sampling & Different randomized \emph{selection algorithm} on AO pairs. \\
\code{rbrp\_id} & Same candidate sampling as \code{rbrp} & Same selection family; additionally changes interpolation/fit via ID. \\
\code{cvt} & No pair sketch; weighted $k$-means++ initialization & Different spatial objective, not random matrix sketching. \\
\bottomrule
\end{longtable}
\normalsize
The three technique rows above use
\code{--anchor-algorithm rand\_sketch}, optionally with
\code{--sketch TECHNIQUE} (default: \code{quadrature}). An explicit sketch
option is invalid for other algorithms. \code{rbrp\_id} is a separate anchor algorithm.

\section{Family A: one factored Gaussian pair map, three inputs}
The implemented sketch entries have the unified form
\begin{equation}
Y_{tg}=\sqrt{\frac{h_g}{\ell}}\,
 (u_t^\mathsf T z_g)(v_t^\mathsf T z_g),
\qquad z_g=f_g\ \text{or}\ a_g/\sqrt{q_g}.
\label{eq:family}
\end{equation}
Equivalently, $Y=S X$ with columns $X_g=\sqrt{h_g}\,(z_g\otimes z_g)$.
Note that $u_t$ and $v_t$ are independent; these are \emph{not} two copies
of the same Gaussian vector. Code generates the factors from two independent
seeded streams. The normalized variant computes the quotient using
$a_g/\sqrt{q_g}$ before projection (algebraically $p_g/q_g$).

\begin{center}
\small
\begin{tabular}{@{}p{0.37\linewidth}p{0.16\linewidth}p{0.33\linewidth}@{}}
\toprule
Sketch technique & $h_g$ & $z_g$ and interpretation \\
\midrule
\code{quadrature} & $|w_g|$ & $f_g$: requested axes, scaled if multiple; AO-only uses raw AO pairs. \\
\code{density} & $m_g$ & $a_g$: raw AO-pair amplitude retained (other axes NYI). \\
\code{density\_col\_norm} & $m_g$ & $a_g/\sqrt{q_g}$: removes pair-column amplitude (other axes NYI). \\
\bottomrule
\end{tabular}
\end{center}
For density modes $m_g$ defaults to spin-summed electron density from a
separate converged RKS B3LYP calculation (\code{density\_weight=electron\_density}).
\code{quadrature\_density} instead sets $m_g=|w_g|\rho_g$; the caller may
supply a complete nonnegative \code{point\_importance} array. There is
\emph{no additional} quadrature factor in either density mode. They require
\code{axes=ao}; zero-weight and zero-importance sites are excluded and sites
with $q_g\leq$ \code{ao\_screening\_threshold} are screened from the density
selection (not from the eventual physical refit). The normalized mode has
$\|X_g\|_2^2=m_g$ because $\|p_g\|=q_g$; raw density instead has
$\|X_g\|_2^2=m_gq_g^2$. This is the meaningful change in its selection
criterion. It is \emph{not} a replacement of the Gaussian Khatri--Rao sketch
by a new sketch distribution, nor a claim of better J/K accuracy.

The three implementations are policies in
\code{cpp/src/anchor\_select\_randomized\_id.cpp}; the general feature
normalization is computed in \code{cpp/src/anchor\_select.cpp}. All three
store $Y\in\mathbb R^{\ell\times n}$, run QRCP on $Y$, and return pivot
\emph{indices only}; they do not currently use the sketch QR factors to fit
interpolation vectors. With explicit rank $k$, $\ell=k+$ oversampling and
$k$ pivots are required if eligible and numerically possible; otherwise
\code{anchor\_sketch\_size} supplies $\ell$, and stopping uses a relative
\emph{sketch} Frobenius-tail tolerance. These tails are not physical pair-space
or Coulomb error guarantees. \code{anchor\_audit} compares physical pair
projection errors to direct QRCP only on small problems.

\section{Family B: random selection without a random sketch matrix}
\textbf{\code{rbrp}, \code{rbrp\_id}.} For AO-only selection, stream symmetric
pair rows
$X_{g,(\mu,\nu)}=\sqrt{|w_g|}a_{g\mu}a_{g\nu}$ with a factor $\sqrt2$
on off-diagonals. This preserves the norm of the ordered-pair representation
without storing duplicate off-diagonal entries. At each iteration, sample a
block of distinct candidate sites with probabilities proportional to their
current squared pair-space residuals; project against the retained subspace,
then run QRCP on the small block and update/refine the residual estimates.
This is an \emph{adaptive residual-weighted sampling algorithm}, not
$S X$ with a fixed Gaussian $S$. The default \code{rbrp} hands the selected
anchors to the usual physical-grid NumPy interpolation fit. \code{rbrp\_id}
selects anchors by the same procedure but constructs ID coefficients and a
pair transform from its factorization (with SVD truncation); that is a
\emph{fit change}, not a new sketch. CPU implementation:
\code{cpp/src/anchor\_select\_rbrp.cpp}; CUDA variant:
\code{cpp/src/anchor\_select\_rbrp\_device.cpp}. Its rank is a cap and the
stopping residual is a physical weighted pair-space projection residual.

\textbf{\code{cvt}.} Importance-weighted centroidal Voronoi clustering uses
random weighted $k$-means++ initialization, minimizes a spatial squared-distance
objective, and snaps centroids to distinct eligible grid sites. Default
importance is density; \code{cvt\_weight=quadrature\_weights} uses $|w_g|$
\emph{instead}, not multiplied by density. It needs a positive rank; QR/sketch
tolerances do not apply. This is a spatial clustering algorithm rather than
pair-space sketching or randomized linear algebra. Implementation:
\code{cpp/src/anchor\_select\_cvt.cpp}.

\section{Family C: deterministic, unsketched pair selection}
\textbf{\code{qrcp\_dense}.} Build $A_g=\sqrt{|w_g|}(f_g\otimes f_g)$ as columns of
an explicit $d^2\times n$ matrix and select grid pivots with LAPACK++ QRCP on
host or MAGMA QRCP on device. A rank is a cap; tolerance compares the QR
diagonal against its first entry. This is the reference selector but its pair
matrix requires $O(d^2n)$ storage. Implementation:
\code{cpp/src/anchor\_select\_qrcp.cpp}.

\textbf{\code{qrcp\_online}.} Evaluate the same feature-pair columns in bounded
batches, incrementally maintaining Householder reflectors and residual norms,
with candidate verification and refreshes. It changes the \emph{execution
algorithm and memory behavior}, not the pair objective and not the random
embedding (there is none). Near pivot ties or numerical thresholds its pivots
need not match dense QR bit for bit. Implementation:
\code{cpp/src/anchor\_select\_online\_qr.cpp}.

\section{Interpreting comparisons}
Changing $|w_g|$ to $m_g$, normalizing $p_g$ by $q_g$, or changing $f_g$ to
$a_g$ changes the \emph{input/objective} seen by QRCP. Moving from dense to
online changes the \emph{algorithm for approximately the same direct
selection task}. Moving to RBRP changes the \emph{randomized selection
mechanism} without defining a Gaussian sketch. Moving to CVT changes the
\emph{objective space} from orbital pairs to real-space geometry. Only
\code{rbrp\_id} replaces the usual post-selection interpolation fit.

For fair comparisons, fix the grid, eligible sites, rank, and relevant
importance convention; use several seeds for randomized methods and inspect
pair-space residuals and downstream J/K accuracy independently. Signed grid
weights enter selection through absolute values; the downstream physical fit
is separate. For usage and caveats see the repository's anchor-selection guide
and README.
\end{document}
